\chapter{A local anomaly and its nonlocal primitive}
\label{ch:local-anomaly}

A local change of an action depends on finitely many derivatives
of its fields at one point. A Green operator instead combines
field values at different points. This difference can decide
whether a symmetry variation has a primitive.

We study a specified cubic variation of an abelian background on
$\mathbb C^2$. Its local cohomology class is nonzero, even when
local coefficients depend arbitrarily and smoothly on position.
A normalized inverse of the real Laplacian gives an explicit
continuous nonlocal primitive. The obstruction to a local
primitive is visible in a polynomial involving two derivative
symbols. A six-generator exterior complex then controls every
possible finite-order local correction.

\section{Background forms and their differential}
\label{sec:la-background}

Write $X=\mathbb C^2=\mathbb R^4$, with
\[
 z_j=x_j+iy_j,\qquad
 \partial_j=\partial_{z_j}
   =\tfrac12(\partial_{x_j}-i\partial_{y_j}),\qquad
 \bar\partial_j=\partial_{\bar z_j}
   =\tfrac12(\partial_{x_j}+i\partial_{y_j}).
\]
Fix the forms and orientation
\begin{equation}\label{eq:la-volume}
 \Omega=dz_1\wedge dz_2,\qquad
 dV_4=dx_1\wedge dy_1\wedge dx_2\wedge dy_2,\qquad
 \nu=d\bar z_1\wedge d\bar z_2\wedge\Omega=4dV_4.
\end{equation}
The last identity follows by expanding $dz_j=dx_j+i\,dy_j$
and $d\bar z_j=dx_j-i\,dy_j$, with the displayed wedge order.
All holomorphic volume forms will remain on the right.

A smooth $(0,q)$-form is a sum of smooth scalar coefficients
times $q$ distinct symbols among $d\bar z_1,d\bar z_2$.
The symbols anticommute. The space of compactly supported such
forms is $\Omega_c^{0,q}(X)$, and
\[
 \mathfrak g_c=\bigoplus_{q=0}^2\Omega_c^{0,q}(X),\qquad
 \bar\partial=\sum_{j=1}^2d\bar z_j\wedge\bar\partial_j.
\]
The scalar derivatives commute, so $\bar\partial^2=0$.
Give this graded vector space the zero Lie bracket. It is the
abelian differential graded Lie algebra used throughout.

To record several inputs and their degrees, use a universal
background
\begin{equation}\label{eq:la-universal}
 A=\sum_{a=1}^m t_a\alpha_a,
 \qquad \alpha_a\in\Omega_c^{0,p_a}(X),
 \qquad |t_a|=1-p_a.
\end{equation}
The parameters generate a free graded commutative algebra over
$\mathbb C$: exchanging degrees $r,s$ gives $(-1)^{rs}$.
In particular, an odd parameter squares to zero. The even
parameters are formal. Every polynomial identity means equality
of coefficients for every finite family in
\eqref{eq:la-universal}. This convention does not require a
basis of the infinite-dimensional space of smooth forms.

The integer degree of a parameter or component coordinate is
called its ghost number. Total degree adds the antiholomorphic
form degree to that number. Put ghost coefficients before forms:
\begin{equation}\label{eq:la-totalization}
 (u\alpha)(v\eta)
 =(-1)^{|\alpha|\operatorname{gh}(v)}
   uv\,\alpha\wedge\eta.
\end{equation}
Here $|\alpha|$ is form degree. The background $A$ has total
degree one. It can be written
\begin{equation}\label{eq:la-components}
 A=c+b_1d\bar z_1+b_2d\bar z_2
       +h\,d\bar z_1\wedge d\bar z_2,
 \qquad |c|=1,\quad |b_j|=0,\quad |h|=-1.
\end{equation}
Thus $c,h$ are odd and $b_1,b_2$ are even. Their scalar
coefficients, at each parameter monomial, are smooth and compactly
supported. The universal odd element satisfies $A^2=0$ over
$\mathbb C$.

The totalized spatial differential is denoted by $d$:
\[
 d(u\alpha)=(-1)^{\operatorname{gh}(u)}
             u\,\bar\partial\alpha.
\]
For a coordinate-dependent coefficient, this rule includes its
spatial derivatives. The degree-one background differential
$\delta$ is a left graded derivation specified by
\begin{equation}\label{eq:la-delta}
 \delta c=0,\qquad
 \delta b_j=-\bar\partial_jc,\qquad
 \delta h=\bar\partial_1b_2-\bar\partial_2b_1.
\end{equation}
It commutes with all scalar spatial derivatives. The two terms
in $\delta^2h$ cancel, and its square on the other components
vanishes directly. Thus $\delta^2=0$ on every polynomial.
In totalized notation,
\begin{equation}\label{eq:la-super-delta}
 \delta A=dA,\qquad \delta d=-d\delta.
\end{equation}
The second sign comes from the increase of ghost number under
$\delta$. Together, \eqref{eq:la-delta} and
\eqref{eq:la-totalization} fix every sign used below.

\section{The cubic variation and a concrete component}
\label{sec:la-cubic}

Set $A_j=\partial_jA$ and define
\begin{equation}\label{eq:la-anomaly}
 B(A)=A_1A_2-A_2A_1,\qquad
 \Theta(A)=-\frac1{24\pi^2}
       \int_X A\,B(A)\wedge\Omega.
\end{equation}
Integration selects antiholomorphic degree two. The three
background factors have total degree three, so $\Theta$ has
ghost number one. Equation~\eqref{eq:la-anomaly} is the given
cubic functional whose primitives we will determine.

There is also a holomorphic exterior notation. Give $dz_j$
degree one in the product rule and put
$\partial A=dz_1\,\partial_1A+dz_2\,\partial_2A$.
Each $\partial_jA$ is odd. Moving the second holomorphic
one-form across it gives
\[
 A\,\partial A\,\partial A
 =-A(\partial_1A\,\partial_2A-\partial_2A\,\partial_1A)
       \wedge\Omega.
\]
Thus \eqref{eq:la-anomaly} is
$+(24\pi^2)^{-1}\int_X A\,\partial A\,\partial A$
in that exterior notation. We retain the even scalar
derivatives $\partial_j$ and the rightmost $\Omega$ below.

A cochain is closed, or a cocycle, when its $\delta$-differential
is zero. It is exact when it equals $\delta C$ for a cochain
$C$ of ghost number one less. Its cohomology class is its class
among cocycles modulo exact cochains. The admissible cochain
space is part of each of these definitions.

\begin{proposition}\label{prop:la-closed}
The functional $\Theta$ is a cubic cocycle of ghost number one.
\end{proposition}
\begin{proof}
The operators $d$ and $\delta$ commute with each $\partial_j$.
The graded product rule gives $\delta B(A)=dB(A)$. Since $A$
is odd,
\[
 \delta\bigl(A B(A)\bigr)
   =(dA)B(A)-A\,dB(A)=d\bigl(A B(A)\bigr).
\]
Each coefficient of this last expression is a spatial derivative
of a compactly supported function. Its integral is zero.
For real derivatives this follows from the fundamental theorem
of calculus and integration in the other coordinates.
The complex derivatives are linear combinations of those real
derivatives. This also justifies the integrations by parts below.
\end{proof}

The following component will distinguish an exact local cochain
from the cocycle in \eqref{eq:la-anomaly}. Take compact smooth
functions $f,g,h$, and ordered odd parameters
$\theta_1,\theta_2,\theta_3$ of degrees $1,1,-1$.
Define $\Theta_3(f,g,h\,d\bar z_1d\bar z_2)$ to be the coefficient
of $\theta_1\theta_2\theta_3$ after setting
\[
 A=\theta_1f+\theta_2g
          +\theta_3h\,d\bar z_1d\bar z_2.
\]
The coefficient convention includes every ordering, with no
factor $1/3!$.

\begin{lemma}\label{lem:la-polarization}
With $f_j=\partial_jf$ and the same notation for $g,h$,
\begin{equation}\label{eq:la-polarized}
 \Theta_3(f,g,h\,d\bar z_1d\bar z_2)
 =-\frac1{4\pi^2}
       \int_X h(f_1g_2-f_2g_1)\,\nu.
\end{equation}
\end{lemma}
\begin{proof}
Each $A_j$ is odd, so $B(A)=2A_1A_2$. The coefficient in
$A B(A)$, with its remaining top form suppressed, is
\[
 2\bigl(fg_1h_2-fh_1g_2-gf_1h_2
          +gh_1f_2+hf_1g_2-hg_1f_2\bigr).
\]
Move derivatives off $h$ in the first four terms. Their integral
becomes
\[
 \int_X h\bigl[-f_2g_1-fg_{12}+f_1g_2+fg_{12}
             +g_2f_1+gf_{12}-g_1f_2-gf_{12}\bigr]\nu.
\]
The mixed second derivatives cancel. The two remaining terms
already contain $h$ without derivatives. All six terms together
therefore give $3h(f_1g_2-f_2g_1)$ inside the factor two.
Multiplication by $-1/(24\pi^2)$ proves the formula.
\end{proof}

For a numerical example, let $\psi$ be a nonnegative smooth
function supported in a ball, with $\int\psi\,dV_4=1$.
Choose a compact smooth $\chi$ equal to one near that ball, and
put $f=\chi z_1$, $g=\chi z_2$, and $h=\psi/4$.
Then the integral in \eqref{eq:la-polarized} equals one, and the
value is $-1/(4\pi^2)$. This proves that $\Theta$ is a nonzero
functional. To prove nonexactness, we must also determine what
every possible local boundary can contain.

\section{All finite-order local cochains}
\label{sec:la-jets}

For $I,J\in\mathbb N^2$, write
$\partial_z^I\partial_{\bar z}^J$
for the indicated product of scalar derivatives. Its order is
$|I|+|J|$, where $|I|=I_1+I_2$.
Introduce jet coordinates
\[
 c_{I,J},\quad (b_1)_{I,J},\quad (b_2)_{I,J},\quad h_{I,J},
\]
representing these derivatives of the four components.
They retain the ghost numbers in \eqref{eq:la-components}.
The algebra of constant-coefficient polynomial densities is
\begin{equation}\label{eq:la-jet-algebra}
 \begin{split}
 \mathcal J={}&
 \operatorname{Sym}\left(
    \bigoplus_{I,J}\mathbb C(b_1)_{I,J}
          \oplus\mathbb C(b_2)_{I,J}\right)\\
 &\otimes
 \Lambda\left(\bigoplus_{I,J}\mathbb Cc_{I,J}
                      \oplus\mathbb Ch_{I,J}\right).
 \end{split}
\end{equation}
The symbol $\operatorname{Sym}$ denotes a polynomial algebra on
the even generators. The exterior algebra $\Lambda$ imposes
anticommutation and square zero on the odd generators.
Each element uses finitely many generators and monomials.

Four even total derivatives act on this algebra. For example,
$D_{z_j}c_{I,J}=c_{I+e_j,J}$ and
$D_{\bar z_j}c_{I,J}=c_{I,J+e_j}$, where $e_j$ is the $j$th
unit multi-index. They act identically on the other component
labels and extend by the Leibniz rule. The differential is
\begin{equation}\label{eq:la-jet-differential}
 \begin{aligned}
 \delta c_{I,J}&=0,&
 \delta(b_j)_{I,J}&=-c_{I,J+e_j},\\
 \delta h_{I,J}&=(b_2)_{I,J+e_1}-(b_1)_{I,J+e_2}.
 \end{aligned}
\end{equation}
It commutes with each total derivative.

Polynomial order counts background factors, including
differentiated factors. Let $\mathcal J_n$ be its order-$n$
part, for $n\geq1$. Define
\begin{equation}\label{eq:la-local-tr}
 \mathcal C_{\mathrm{loc,tr},n}
 =\mathcal J_n\Big/
  \left(\sum_{j=1}^2D_{z_j}\mathcal J_n
             +\sum_{j=1}^2D_{\bar z_j}\mathcal J_n\right).
\end{equation}
This is a vector-space quotient by sums of total derivatives.
It is not a quotient by an ideal of densities. Evaluation is
\[
 [P]\longmapsto \int_X P(jA)\,\nu,
\]
where $jA$ denotes the collection of component jets.
Every total derivative has zero integral on compact inputs.
Since $\delta$ preserves the displayed subspace, it induces a
differential on the quotient.

All complex multilinear maps of form components occur in
\eqref{eq:la-jet-algebra}, after polarization. For instance,
ghost-zero cubic densities may use three $b$ components, or one
component of each type $c,b,h$, with any derivatives.
No form-wedge restriction is imposed on these densities.
Both holomorphic and antiholomorphic jets are included.

Allowing coefficients to depend on position gives the algebra
of finite sums
\begin{equation}\label{eq:la-smooth-jets}
 \mathcal J_{\mathrm{sm}}
 =\left\{\sum_\mu f_\mu(z,\bar z)M_\mu:
       f_\mu\in C^\infty(X),\quad M_\mu\in\mathcal J
       \text{ a monomial}\right\}.
\end{equation}
Set $\delta f_\mu=0$. Total derivatives now differentiate the
coefficients as well as the jets. Define
\begin{equation}\label{eq:la-local-sm}
 \mathcal C_{\mathrm{loc,sm},n}
 = (\mathcal J_{\mathrm{sm}})_n\Big/
  \left(\sum_jD_{z_j}(\mathcal J_{\mathrm{sm}})_n
           +\sum_jD_{\bar z_j}(\mathcal J_{\mathrm{sm}})_n\right).
\end{equation}
Compact input support makes each integral finite, regardless of
the growth of the smooth coefficients. A locally bounded
finite-order density means an expression of this kind on each
relatively compact coordinate neighborhood, with compatible
expressions where neighborhoods overlap.

The differential preserves polynomial order. A primitive for
$\Theta$ would therefore have a cubic ghost-zero component
that is itself a primitive. We can restrict the proof to $n=3$.
The same observation applies to the formal product over all
polynomial orders, with finite jet order at each fixed order.

\section{Polarization and integration by parts}
\label{sec:la-normal-form}

Polarization replaces a cubic expression by three labeled
inputs and extracts the coefficient in which each input occurs
once. Exchanging homogeneous component labels of degrees $r,s$
gives the sign $(-1)^{rs}$. These are the tensor signs used here.
The inverse construction identifies the three inputs and divides
by $3!$. Both operations are valid over $\mathbb C$.
For odd inputs they are interpreted using independent graded
parameters, as in Lemma~\ref{lem:la-polarization}.

It is useful first to keep the three labels ordered. Their
permutations form the group $S_3$. The averaging operator
\[
 \operatorname{Av}(v)=\frac16\sum_{\sigma\in S_3}\sigma v
\]
projects onto the graded invariant tensors. It commutes with
$\delta$. It also identifies invariants with the quotient by
vectors $v-\sigma v$: averaging vanishes on those differences,
and its class in that quotient is the class of $v$.

Let $\nabla=(\partial_1,\partial_2,
\bar\partial_1,\bar\partial_2)$, and use four-component
multi-indices for its powers in this section. Fix one component
label in each slot. An ordered polarized density has the form
\[
 \int_X\sum_{I,J,K}f_{I,J,K}(x)
       (\nabla^I\phi_1)(\nabla^J\phi_2)
                         (\nabla^K\phi_3)\,\nu.
\]
Every sum is finite. Iterated integration by parts puts it in
the normal form
\begin{equation}\label{eq:la-normal-form}
 \int_X\sum_{I,J}g_{I,J}(x)
       (\nabla^I\phi_1)(\nabla^J\phi_2)\phi_3\,\nu.
\end{equation}
For one derivative on the last slot, the rule is
\[
 \int_X fU V(\nabla_a\phi_3)\nu
 =-\int_X\bigl[(\nabla_af)UV
               +f(\nabla_aU)V+fU(\nabla_aV)\bigr]\phi_3\nu.
\]
Spatial derivatives are even, so this step creates no additional
ghost sign.

\begin{lemma}\label{lem:la-normal-unique}
For ordered inputs and fixed component labels, normal form
\eqref{eq:la-normal-form} is unique. It identifies the density
quotient with its space of multilinear functionals on compact
smooth inputs.
\end{lemma}
\begin{proof}
Suppose a normal form evaluates to zero on every three inputs.
Fix $\phi_1,\phi_2$. Varying the compact test $\phi_3$ shows that
the smooth function
\[
 \sum_{I,J}g_{I,J}(x)
          (\nabla^I\phi_1)(x)(\nabla^J\phi_2)(x)
\]
vanishes pointwise. To justify this implication, a nonzero
continuous value has a small neighborhood on which an
appropriate real part has one sign. A nonnegative bump
function there would give a nonzero integral.

At any point $x$, independently prescribe the finitely many
jets of $\phi_1$ and $\phi_2$. A Taylor polynomial multiplied
by a bump function equal to one near $x$ realizes these jets.
The real and complex derivative bases are related by an
invertible linear change, so either basis permits arbitrary
finite jets. Taking one jet at a time proves
$g_{I,J}(x)=0$ for every pair. This applies independently to
all component labels. Odd coefficients are tested with distinct
formal parameters, so their signs impose no further relation
between ordered slots.

Every reduction to normal form was a sum of integrations by
parts. A zero functional has zero normal form, and therefore
already represents zero in the displayed density quotient.
\end{proof}

For constant coefficients, assign symbols $p_{ij},q_{ij}$ to
$\partial_j,\bar\partial_j$ in slot $i$. They are commuting
polynomial variables. Total derivatives act by multiplication
by the sums of the corresponding slot symbols. Thus the ordered
symbol ring is
\begin{equation}\label{eq:la-symbol-ring}
 R=\mathbb C[p_{ij},q_{ij}:1\leq i\leq3,\ 1\leq j\leq2]
       \Big/\left(\sum_i p_{ij},\sum_iq_{ij}:j=1,2\right).
\end{equation}
Here the module quotient by total momentum is exactly the
polarized realization of the vector-space quotient
\eqref{eq:la-local-tr}. Normal-form uniqueness proves that it
introduces all and only the integration-by-parts relations.

\section{The nonzero local symbol}
\label{sec:la-symbol-obstruction}

Each ordered slot has component labels $c_i,b_{i1},b_{i2},h_i$.
Equations~\eqref{eq:la-jet-differential} become
\begin{equation}\label{eq:la-symbol-differential}
 \delta c_i=0,\qquad
 \delta b_{ij}=-q_{ij}c_i,\qquad
 \delta h_i=q_{i1}b_{i2}-q_{i2}b_{i1}.
\end{equation}
On a tensor product, $\delta$ acting in slot $i$ has the sign
given by the sum of the preceding component degrees.

Set all antiholomorphic symbols to zero:
\begin{equation}\label{eq:la-qzero}
 q_{ij}=0\qquad(1\leq i\leq3,\ 1\leq j\leq2).
\end{equation}
Every differential in \eqref{eq:la-symbol-differential} then
vanishes. Therefore specialization \eqref{eq:la-qzero} sends
every boundary to zero. It respects all relations in
\eqref{eq:la-symbol-ring} and all permutations of slots.

Lemma~\ref{lem:la-polarization} gives the following component of
the specialized anomaly, with input form degrees $(0,0,2)$:
\begin{equation}\label{eq:la-determinant-symbol}
 -\frac1{4\pi^2}(p_{11}p_{22}-p_{12}p_{21}).
\end{equation}
This polynomial is nonzero under the holomorphic momentum
relations. Indeed, the values
\[
 (p_{11},p_{12})=(1,0),\qquad
 (p_{21},p_{22})=(0,1),\qquad
 (p_{31},p_{32})=(-1,-1)
\]
satisfy those relations and give $-1/(4\pi^2)$.
They are values in the complex polynomial ring of derivative
symbols. No compactly supported holomorphic input is required
for this algebraic test.

\begin{proposition}\label{prop:la-tr-nonexact}
The class of $\Theta$ is nonzero in
$H^1(\mathcal C_{\mathrm{loc,tr},3},\delta)$.
\end{proposition}
\begin{proof}
Proposition~\ref{prop:la-closed} proves closedness.
A boundary would vanish under \eqref{eq:la-qzero}, whereas
\eqref{eq:la-determinant-symbol} does not. The component is
taken after graded polarization, so this contradiction holds
in the invariant subspace as well as in the ordered complex.
\end{proof}

With position-dependent coefficients, integration by parts
can differentiate the coefficient and remove a background
derivative. The next calculation determines whether such
lower-order terms can supply a primitive.

\section{The complete component complex}
\label{sec:la-koszul}

An exterior algebra gives a convenient description of all
component labels. For slot $i$, introduce a one-dimensional
symbol $\kappa_i$ of degree one and exterior generators
$e_{i1},e_{i2}$ of degree minus one. Identify
\begin{equation}\label{eq:la-slot-identification}
 c_i\longleftrightarrow\kappa_i\otimes1,\qquad
 b_{ij}\longleftrightarrow\kappa_i\otimes e_{ij},\qquad
 h_i\longleftrightarrow-\kappa_i\otimes e_{i1}e_{i2}.
\end{equation}
For exterior generators, contraction $\iota_q$ is the odd
derivation defined by $\iota_q(e_{ij})=q_{ij}$.
For example,
$\iota_q(e_{i1}e_{i2})=q_{i1}e_{i2}-q_{i2}e_{i1}$.
Set
\[
 d(\kappa_i\otimes\omega)
       =-\kappa_i\otimes\iota_q\omega.
\]
This gives $d b_{ij}=-q_{ij}c_i$ and
$d h_i=q_{i1}b_{i2}-q_{i2}b_{i1}$, including the sign of the
two-form component.

In the tensor product of three slots, move the $\kappa_i$'s to
the front. If the exterior degrees in the slots are $k_i$, the
permutation sign is $(-1)^{2k_1+k_2}$. Combining it with the
tensor differential gives the full complex
\begin{equation}\label{eq:la-koszul-complex}
 K=R\otimes\Lambda(e_{11},e_{12},e_{21},e_{22},e_{31},e_{32}),
 \qquad d_K=-\iota_q.
\end{equation}
The suppressed factor $\kappa_1\kappa_2\kappa_3$ has degree
three. Hence exterior degree $k$ in \eqref{eq:la-koszul-complex}
has ghost degree $3-k$. The operator $d_K$ raises ghost degree
by one. Its square vanishes because contractions of distinct
exterior generators anticommute, while their coefficients
$q_{ij}$ commute. This exterior complex with differential given
by contraction is called a Koszul complex.

The momentum relations give $q_{3j}=-q_{1j}-q_{2j}$.
Change the exterior basis to
\begin{equation}\label{eq:la-redundant-generators}
 e_{11},e_{12},e_{21},e_{22},\qquad
 r_j=e_{1j}+e_{2j}+e_{3j}\quad(j=1,2).
\end{equation}
The change is invertible, since
$e_{3j}=r_j-e_{1j}-e_{2j}$. The last two generators have zero
differential. Call the first four exterior generators active,
and enumerate their momenta as $q_a$, $1\leq a\leq4$.
The coefficient ring is
\[
 R=\mathbb C[p_{11},p_{12},p_{21},p_{22},
             q_{11},q_{12},q_{21},q_{22}].
\]

Define an operator and two counting operators by
\begin{equation}\label{eq:la-koszul-homotopy}
 h_K=-\sum_{a=1}^4e_a\frac{\partial}{\partial q_a},
 \qquad
 N_q=\sum_aq_a\frac{\partial}{\partial q_a},
 \qquad N_e=\sum_ae_a\iota_a.
\end{equation}
Here $\iota_a$ deletes $e_a$ with the sign of its position in
an exterior monomial. It annihilates the other basis generators,
including $r_1,r_2$. In particular,
\[
 \iota_a e_b+e_b\iota_a=\delta_{ab},\qquad
 \frac{\partial}{\partial q_b}q_a
       =q_a\frac{\partial}{\partial q_b}+\delta_{ab}.
\]
Expanding both compositions using these two identities gives
\begin{equation}\label{eq:la-euler-identity}
 d_Kh_K+h_Kd_K=N_q+N_e.
\end{equation}
More explicitly, $d_Kh_K$ contains
$\sum_aq_a\partial_{q_a}
 -\sum_{a,b}q_ae_b\iota_a\partial_{q_b}$.
The other composition contains
$\sum_ae_a\iota_a
 +\sum_{a,b}e_bq_a\partial_{q_b}\iota_a$.
The double sums cancel.

The nonnegative integer $N=N_q+N_e$ counts active momentum
degree plus active exterior degree. Both $d_K$ and $h_K$
preserve $N$. Every polynomial is a finite sum of its
$N$-homogeneous parts. For $N>0$, division by this integer in
\eqref{eq:la-euler-identity} gives a contraction $h_K/N$.
The subspace $N=0$ consists exactly of
\begin{equation}\label{eq:la-koszul-cohomology}
 \mathbb C[p_{11},p_{12},p_{21},p_{22}]
                       \otimes\Lambda(r_1,r_2).
\end{equation}
Its differential is zero, and it cannot receive a boundary
from a positive-$N$ part. Thus \eqref{eq:la-koszul-cohomology}
is the cohomology of the full ordered complex. Its ghost
degrees are $3,2,1$.

\begin{lemma}\label{lem:la-order-lowering}
The cubic constant-coefficient symbol complex has zero
cohomology in ghost degree zero. A closed homogeneous symbol
of ghost degree zero and differential order $r$ has a primitive
of order $r-1$. The same statement holds for graded invariant
symbols and for smooth families of symbols.
\end{lemma}
\begin{proof}
Ghost degree zero requires exterior degree three. There are
only two generators $r_1,r_2$, so at least one active exterior
generator occurs. Every such component has $N>0$.
For a closed component $v$, equation~\eqref{eq:la-euler-identity}
therefore gives
\[
 d_K\left(\frac{h_Kv}{N}\right)=v.
\]
The operator $h_K$ differentiates once in an antiholomorphic
momentum. It lowers total momentum order by one and lowers
ghost degree by one. Holomorphic momentum factors remain
unrestricted. Sum over the finitely many $N$-components.
If $r=0$, the operator has zero output, so closedness forces
$v=0$.

For an invariant $v$, average a primitive over $S_3$.
Since permutations commute with $d_K$, its differential remains
$v$. Averaging preserves polynomial order. All operations are
linear on the finite coefficient array of a polynomial.
Applying them to smooth coefficient functions therefore
preserves their smooth dependence.
\end{proof}

\section{Position-dependent local counterterms}
\label{sec:la-smooth-obstruction}

Filter normal form \eqref{eq:la-normal-form} by
$|I|+|J|$, the total derivative order on its background slots.
Derivatives of the coefficient functions are not counted.
The top homogeneous derivative polynomial is its principal
symbol. The operator $\delta$ raises the filtration by at most
one. Its principal differential is
\eqref{eq:la-symbol-differential} with $q_3=-q_1-q_2$.
In fact, a derivative introduced on the third slot moves onto
the first two slots with this minus sign. A term in which it
instead differentiates a coefficient has lower order.

Permuting slots preserves the filtration: integration by parts
does not increase the total number of background derivatives.
Its top action is the permutation action in the symbol ring.
Consequently a symmetric top symbol can be lifted and then
averaged without changing that symbol. This permits applying
Lemma~\ref{lem:la-order-lowering} to actual cubic functionals.

\begin{theorem}\label{thm:la-smooth-nonexact}
The cocycle $\Theta$ is not exact in
$\mathcal C_{\mathrm{loc,sm},3}$.
It also has no primitive with locally bounded finite jet order
and arbitrary smooth coefficients.
\end{theorem}
\begin{proof}
Suppose $\delta C=\Theta$, where $C$ has ghost degree zero and
finite derivative order $r$. If $r\geq2$, the order-$(r+1)$
symbol of $\delta C$ is zero, since $\Theta$ has order two.
Thus the top symbol $\sigma_r(C)$ is a closed ghost-zero
symbol.

Lemma~\ref{lem:la-order-lowering}, applied to its smooth
coefficients, gives a ghost-minus-one symbol $E$ of order
$r-1$ such that $d_KE=\sigma_r(C)$. Lift $E$ to normal form
and average it to obtain a cubic functional $\widetilde E$.
Then
\[
 C\longmapsto C-\delta\widetilde E
\]
preserves the equation $\delta C=\Theta$ and strictly lowers
the order of $C$. Coefficient derivatives contribute only at
that lower order. Repetition terminates because $r$ is a
nonnegative integer.

The resulting primitive would have order at most one. If its
order is one, the order-two equation makes the symbol of
$\Theta$ a boundary in the constant-coefficient symbol
complex, at every point $x$. Specialization
\eqref{eq:la-qzero} annihilates that boundary. Its component
\eqref{eq:la-determinant-symbol} remains nonzero, a
contradiction. If the order is zero, its differential has
order at most one and cannot equal the nonzero order-two
symbol of $\Theta$.

All arguments hold on a nonempty coordinate ball, using
inputs compactly supported inside it. A locally bounded
finite-order primitive restricts to a finite-order primitive
on a smaller such ball. The same contradiction excludes it.
\end{proof}

The order reduction is needed because a smooth coefficient can
produce lower-order terms under integration by parts. It proves
nonexactness in the full smooth local complex, including all
those terms. An unrestricted inverse differential operator
requires a different cochain space, which we now define.

\section{Continuous cochains on compact smooth inputs}
\label{sec:la-continuous}

For a compact set $K\subset X$, forms supported in $K$ have
seminorms
\[
 \|\alpha\|_{K,m}
 =\sum_{|\mu|\leq m}\sup_{x\in K}
                      |\partial_{\mathbb R}^{\mu}\alpha(x)|,
 \qquad m\geq0.
\]
The derivatives are real coordinate derivatives, and the norm
uses the four form coefficients. Their resulting topology is
the fixed-support smooth topology. Choose an exhaustion by
compact balls. The usual topology on $\mathfrak g_c$ is the
locally convex inductive limit of these spaces, denoted the
LF topology. It is the finest locally convex topology making
all fixed-support inclusions continuous.

In particular, a seminorm whose restriction to each
fixed-support space is continuous is continuous in this
topology. To see this directly, that seminorm defines a
locally convex topology making every inclusion continuous.
The defining finest topology is at least as strong.

For an explicit cochain convention, let $s\alpha$ have degree
$q-1$ when $\alpha\in\Omega_c^{0,q}(X)$.
This is the underlying graded shift $\mathfrak g_c[1]$ with
the convention $E[k]^j=E^{j+k}$. On its graded symmetric
tensors, prescribe the degree-one operator
\begin{equation}\label{eq:la-cochain-input-differential}
 q_1(s\alpha)=s\bar\partial\alpha,
 \qquad
 q_1(x_1\cdots x_n)
 =\sum_i(-1)^{\sum_{j<i}|x_j|}
       x_1\cdots(q_1x_i)\cdots x_n.
\end{equation}
The sign in the first formula is part of the convention.
The relation $q_1^2=0$ follows from $\bar\partial^2=0$ and
the tensor signs.

A degree-$k$ continuous cochain of polynomial order $n$ is a
jointly continuous graded symmetric multilinear functional
on these shifted inputs, with scalar values. It vanishes
unless $\sum_i(q_i-1)=-k$. Define
\begin{equation}\label{eq:la-continuous-differential}
 \delta F=-(-1)^kFq_1.
\end{equation}
This differential has degree one and squares to zero.
On polynomial coordinates it is a left derivation. For a
component coordinate of ghost number $1-q$, the prefactor in
\eqref{eq:la-continuous-differential} is $(-1)^q$.
It therefore reproduces exactly \eqref{eq:la-delta}.
Equivalently, evaluate every formula on the finite backgrounds
\eqref{eq:la-universal} and apply those component rules.

Denote this complex by $\mathcal C_{\mathrm{cont},n}$.
All differentiations of inputs remain continuous, since they
raise a fixed-support derivative seminorm by one and preserve
support. In polynomial order three, local finite-order
cochains with smooth coefficients define continuous cochains
in this sense. For a finite sum with coefficients $f_\mu$,
the smooth positive weight $w=1+\sum_\mu|f_\mu|^2$
dominates every $|f_\mu|$. Weighted supremum seminorms of
the first input, ordinary derivative suprema of the second,
and derivative $L^1$ seminorms of the third bound each integral.
Each of these seminorms is continuous by the LF criterion
above. This also proves joint continuity rather than only
continuity for a common fixed support.

\section{A normalized Green homotopy}
\label{sec:la-green}

Let $r=|z|$ and define
\begin{equation}\label{eq:la-green-function}
 G_4(z)=\frac1{4\pi^2r^2}\qquad(r>0).
\end{equation}
It is locally integrable, since the radial volume factor in
four dimensions is $r^3dr$. For a smooth radial function,
direct differentiation gives
\[
 \Delta_{\mathbb R^4}g(r)=g''(r)+3r^{-1}g'(r).
\]
Thus $\Delta G_4=0$ away from the origin.

The area of the unit three-sphere is $2\pi^2$. One derivation
fixes the normalization without an orientation choice left
implicit. The one-dimensional Gaussian integral $I$ satisfies
\[
 I^2=\int_{\mathbb R^2}e^{-(x^2+y^2)}dx\,dy
       =2\pi\int_0^\infty re^{-r^2}dr=\pi.
\]
Positivity gives $I=\sqrt\pi$. Taking four factors and using
four-dimensional polar coordinates gives
\[
 \pi^2=\operatorname{area}(S^3)
              \int_0^\infty r^3e^{-r^2}dr
         =\tfrac12\operatorname{area}(S^3).
\]
All integrands are positive, so iterated integration and the
radial change of variables give the displayed equalities.

\begin{lemma}\label{lem:la-green-delta}
As a distribution on $\mathbb R^4$,
\begin{equation}\label{eq:la-green-delta}
 (-\Delta_{\mathbb R^4})G_4=\delta_0,
\end{equation}
where $\delta_0(\phi)=\phi(0)$ for a compact smooth test
function $\phi$.
\end{lemma}
\begin{proof}
A locally integrable function acts on tests by integration.
Distributional differentiation is defined by moving each
derivative onto the test with a minus sign. It is therefore
enough to prove $\int G_4(-\Delta\phi)dV_4=\phi(0)$.

Integrate by parts twice on $\epsilon<r<R$, with $R$ outside
the support of $\phi$. The outer boundary contributes zero.
At the inner sphere the outward normal of this region is
$-\partial_r$. Since $G'_4(r)=-1/(2\pi^2r^3)$, the result is
\[
 \int_{\epsilon<r<R}G_4(-\Delta\phi)dV_4
 =\int_{r=\epsilon}
     \left[\frac{\phi}{2\pi^2\epsilon^3}
                  +G_4\partial_r\phi\right]dS.
\]
The first term is the spherical average of $\phi$. Its
difference from $\phi(0)$ is at most
$\epsilon\sup_{r\leq\epsilon}|\nabla\phi|$.
The second term has absolute value at most
$\epsilon\sup|\nabla\phi|/2$.
Finally,
\[
 \int_{r<\epsilon}G_4\,dV_4=\frac{\epsilon^2}{4},
\]
so the omitted integral against $\Delta\phi$ tends to zero.
Letting $\epsilon\to0$ proves \eqref{eq:la-green-delta}.
\end{proof}

For constant exterior symbols, let $\iota_j$ delete
$d\bar z_j$ from a wedge monomial, with the sign of its
position, and annihilate monomials without that symbol.
It satisfies
\[
 (d\bar z_i\wedge)\iota_j
       +\iota_j(d\bar z_i\wedge)=\delta_{ij}.
\]
Define an operator lowering form degree by one,
\[
 R=-4\sum_{j=1}^2\partial_j\iota_j.
\]
The preceding relation and commuting scalar derivatives give
\begin{equation}\label{eq:la-laplacian-homotopy}
 \bar\partial R+R\bar\partial
 =-4\sum_j\partial_j\bar\partial_j
 =-\Delta_{\mathbb R^4}.
\end{equation}

For a compact smooth form $f$, convolve its coefficients with
$G_4$ using real volume, and set
\begin{equation}\label{eq:la-hraw}
 (G_4*f)(z)=\int_XG_4(z-w)f(w)dV_w,
 \qquad H_{\mathrm{raw}}f=R(G_4*f).
\end{equation}
The output is smooth. To justify differentiation, write the
convolution using $u=z-w$ and move derivatives to $f(z-u)$.
For $z$ in a fixed compact set, all relevant $u$ lie in a
fixed compact set. Local integrability of $G_4$ bounds the
integral of every such derivative. Difference quotients then
converge under that bound. Distributional differentiation
agrees with this computation.

Equations~\eqref{eq:la-green-delta} and
\eqref{eq:la-laplacian-homotopy} now prove
\begin{equation}\label{eq:la-homotopy-raw}
 \bar\partial H_{\mathrm{raw}}
       +H_{\mathrm{raw}}\bar\partial=1
 \quad\text{on }\Omega_c^{0,\bullet}(X).
\end{equation}
This is an identity into all smooth forms. Compact support of
the output is not required.

Differentiating \eqref{eq:la-green-function} gives the explicit
integral
\begin{equation}\label{eq:la-hraw-integral}
 (H_{\mathrm{raw}}f)(z)
 =\frac1{\pi^2}\sum_{j=1}^2
      \int_X\frac{\overline{z_j-w_j}}{|z-w|^4}\,
                    \iota_jf(w)\,dV_w.
\end{equation}
Its kernel has size at most a constant times $|z-w|^{-3}$,
which is locally integrable in four dimensions. The ordinary
first derivative of $G_4$ is its distributional first
derivative: the boundary term on a sphere of radius
$\epsilon$ is bounded by a constant times $\epsilon$ and
tends to zero. Thus \eqref{eq:la-hraw-integral} includes the
entire kernel, with no additional term at the origin.

\section{The nonlocal primitive and its continuity}
\label{sec:la-primitive}

Extend $H_{\mathrm{raw}}$ to formal coefficients by
\begin{equation}\label{eq:la-h-total}
 \mathsf H(u\alpha)
 =(-1)^{\operatorname{gh}(u)}uH_{\mathrm{raw}}\alpha.
\end{equation}
This is the totalization for an odd spatial operator.
Applying it twice in the two compositions of
\eqref{eq:la-homotopy-raw} gives
\begin{equation}\label{eq:la-h-total-identities}
 d\mathsf H+\mathsf Hd=1,
 \qquad \delta\mathsf H=-\mathsf H\delta.
\end{equation}
For the second identity, $\delta$ increases the ghost number
of the coefficient by one, reversing the sign in
\eqref{eq:la-h-total}. Equivalently, the component of
$\mathsf HA$ of form degree $r$ is
$(-1)^rH_{\mathrm{raw}}A_{r+1}$. This describes the sign for
both the one-form and two-form background inputs.

Define
\begin{equation}\label{eq:la-primitive}
 \boxed{\displaystyle
 \Psi(A)=\frac1{24\pi^2}
       \int_X(\mathsf HA)
        (\partial_1A\,\partial_2A
                  -\partial_2A\,\partial_1A)\wedge\Omega.}
\end{equation}
The element $\mathsf HA$ has total degree zero. The other
two factors have total degree two together. Integration
selects form degree two, so $\Psi$ has ghost degree zero.
It is a homogeneous cubic expression.

\begin{proposition}\label{prop:la-primitive-identity}
For every compact smooth universal background,
\begin{equation}\label{eq:la-primitive-identity}
 \delta\Psi=\Theta.
\end{equation}
\end{proposition}
\begin{proof}
Use $B(A)$ from \eqref{eq:la-anomaly}.
Equations~\eqref{eq:la-super-delta} and
\eqref{eq:la-h-total-identities} give
\[
 \delta(\mathsf HA)=-\mathsf HdA,
 \qquad \delta B(A)=dB(A).
\]
Since $\mathsf HA$ is even, the product rule and integration
by parts give
\begin{align*}
 \delta\Psi
 &=\frac1{24\pi^2}\int_X
       \bigl[-(\mathsf HdA)B(A)
                     +(\mathsf HA)dB(A)\bigr]\wedge\Omega\\
 &=-\frac1{24\pi^2}\int_X
       \bigl[\mathsf HdA+d\mathsf HA\bigr]B(A)\wedge\Omega\\
 &=-\frac1{24\pi^2}\int_X A B(A)\wedge\Omega
 =\Theta(A).
\end{align*}
Every coefficient of $B(A)$ has compact support, while
$\mathsf HA$ is smooth. Thus $(\mathsf HA)B(A)$ has compact
support, which justifies the integration by parts. In a
polarized term, the two differentiated inputs have compact
support independently. Their product still has compact
support, regardless of the support of the Green output.
\end{proof}

Continuity must hold when the supports vary. For any compact
smooth form $f$, split the integral
\eqref{eq:la-hraw-integral} into $|z-w|\leq1$ and its
complement. On the first region,
$\int_{|u|\leq1}|u|^{-3}dV_u=2\pi^2$.
On the second region, $|u|^{-3}\leq1$. Hence a constant
independent of both support and $z$ gives
\begin{equation}\label{eq:la-h-bound}
 \|H_{\mathrm{raw}}f\|_\infty
 \leq C\bigl(\|f\|_\infty+\|f\|_{L^1}\bigr).
\end{equation}
Define a global seminorm using the real derivatives:
\begin{equation}\label{eq:la-pone}
 p_1(f)=\sum_{|\mu|\leq1}
       \left(\|\partial_{\mathbb R}^{\mu}f\|_\infty
                  +\|\partial_{\mathbb R}^{\mu}f\|_{L^1}\right).
\end{equation}
Its restriction to a fixed-support space is continuous:
the integral norms are bounded by the support volume times
the corresponding supremum norms. The LF criterion in
Section~\ref{sec:la-continuous} proves continuity on the full
compact-support space.

Each term in the polarization of $\Psi$ is bounded by
\[
 C\|H_{\mathrm{raw}}\alpha\|_\infty
       \|\partial\beta\|_\infty
       \|\partial\gamma\|_{L^1}.
\]
There are finitely many permutations and component terms.
Using \eqref{eq:la-h-bound} and
\eqref{eq:la-pone} therefore gives a support-independent
constant $C'$ such that
\begin{equation}\label{eq:la-joint-bound}
 |\Psi_3(\alpha,\beta,\gamma)|
     \leq C'p_1(\alpha)p_1(\beta)p_1(\gamma).
\end{equation}
This product of continuous seminorms proves joint continuity
at zero. Multilinearity, by expanding increments in each
argument, proves joint continuity at every triple.
Thus $\Psi\in\mathcal C_{\mathrm{cont},3}^0$.

The kernel in \eqref{eq:la-hraw-integral} depends on the
difference $z-w$. Changes of variables prove that $\Psi$
is translation invariant. Its formula uses one input at
$w$ and two differentiated inputs at $z$, connected by this
kernel. All integrals act on smooth inputs; local integrability
handles the single singular difference. Theorem
\ref{thm:la-smooth-nonexact} and
\eqref{eq:la-primitive-identity} also show that this functional
has no finite-order local representation.

\section{Local cancellation and the choice of cochains}
\label{sec:la-consequence}

\begin{theorem}\label{thm:la-three-parts}
For the abelian Dolbeault algebra
$\Omega_c^{0,\bullet}(\mathbb C^2)$ and the cubic functional
$\Theta$ of \eqref{eq:la-anomaly}, the following hold.
\begin{enumerate}
 \item Its class is nonzero among constant-coefficient
 finite-order local polynomial cochains, with all form
 components and all spatial jets allowed.
 \item Its class remains nonzero with arbitrary smooth
 local coefficients, including locally bounded finite jet
 order.
 \item It is exact among jointly continuous cochains on
 compact smooth inputs. The translation-invariant nonlocal
 primitive is \eqref{eq:la-primitive}.
\end{enumerate}
\end{theorem}
\begin{proof}
The first part is Proposition~\ref{prop:la-tr-nonexact}.
The second is Theorem~\ref{thm:la-smooth-nonexact}.
For the third, Proposition~\ref{prop:la-primitive-identity}
gives exactness, and \eqref{eq:la-joint-bound} proves the
required continuity. Polynomial-order preservation extends
the first two statements to arbitrary local polynomial
counterterms and to formal completion by polynomial order.
\end{proof}

In the scalar local complex defined above, a counterterm is
a ghost-zero local functional added to the scalar action.
Its change of the cubic variation is a $\delta$-boundary.
Thus a nonzero multiple of $\Theta$ cannot be canceled by
such a counterterm. A continuous nonlocal change can cancel
it by adding the opposite multiple of $\Psi$.
This consequence concerns the specified scalar background
complex. A theory with additional fields or nonlinear
differentials requires its own counterterm complex and a
comparison with this one.

For $\lambda\in\mathbb C$, the local class of $\lambda\Theta$
is nonzero exactly when $\lambda\ne0$. Multiplication by a
nonzero scalar is invertible on these complexes. When
$\lambda=0$, the functional itself is zero. If independent
charge contributions are defined to be $q_j^3\Theta$, their
sum is $(\sum_jq_j^3)\Theta$. Charges $1,-1$ give zero, while
charges $1,1$ give twice the nonzero local class. These are
statements about the specified cubic functional and its
scalar multiples.

The finite-order condition is imposed at each fixed
polynomial order. A completion admitting infinitely many
derivatives at that same order is a different space and
need not act on arbitrary compact smooth inputs. The Green
primitive also uses the whole Euclidean space and allows
noncompact output. The proofs above make no assertion for
unrestricted noncompact inputs, other boundary conditions,
or a compatible local primitive on every open set. Within
the stated domains, the same cubic variation has a nonzero
local class and the explicit continuous primitive
\eqref{eq:la-primitive}.
